\documentclass[10pt]{extarticle}
\usepackage[a4paper,margin=1.5cm]{geometry}
\usepackage[T1]{fontenc}
\usepackage{amsmath,amssymb}
\usepackage{graphicx}
\usepackage[export]{adjustbox}
\usepackage{multicol}
\usepackage{enumitem}
\usepackage{booktabs,array,tabularx}
\setlength{\columnsep}{0.7cm}
\setlength{\parindent}{0pt}
\setlength{\parskip}{3pt}
\pagestyle{empty}
\makeatletter
\newcommand{\qfigdef}[4]{\@namedef{qf@file@#1}{#2}\@namedef{qf@w@#1}{#3}\@namedef{qf@wide@#1}{#4}}
\newcommand{\qf@check}[1]{\@ifundefined{qf@file@#1}{\PackageError{score2}{Unknown figure id #1}{}}{}}
\newcommand{\qf@img}[4][]{\includegraphics[width=#2\dimexpr\@nameuse{qf@w@#4}pt\relax,max width=#3,max height=0.45\textheight,#1]{\@nameuse{qf@file@#4}}}
\newif\ifqf@full
\newif\ifqf@boxed
\newsavebox{\qf@box}
\newcommand{\qf@open}{\global\qf@fulltrue\par\global\qf@boxedtrue
  \setbox\qf@box\vbox\bgroup\hsize\textwidth\linewidth\textwidth}
\newcommand{\qf@emit}[1]{\par\egroup\global\qf@boxedfalse
  \ifdim\dimexpr\ht\qf@box+\dp\qf@box\relax<#1\textheight\noindent\usebox{\qf@box}\par\else\unvbox\qf@box\fi}
\newcommand{\qfullsplit}{\ifqf@boxed\qf@emit{0.85}\fi}
\newcommand{\qf@close}{\ifqf@boxed\qf@emit{0.45}\fi\global\qf@fullfalse}
\newcommand{\qfullbegin}{\end{multicols}\qf@open}
\newcommand{\qfullend}{\qf@close\begin{multicols}{2}}
\newcommand{\qfullstart}{\qf@open}
\newcommand{\qfulldone}{\qf@close}
\newcommand{\qfig}[2][1]{\qf@check{#2}%
  \ifqf@full\qf@blk{#1}{#2}\else
  \ifnum\@nameuse{qf@wide@#2}=1 \expandafter\@firstoftwo\else\expandafter\@secondoftwo\fi
  {\end{multicols}\begin{center}\qf@img{#1}{\textwidth}{#2}\end{center}\begin{multicols}{2}}%
  {\qf@blk{#1}{#2}}\fi}
\newcommand{\qf@blk}[2]{\par\smallskip{\centering\qf@img{#1}{\linewidth}{#2}\par}\smallskip}
\newcommand{\qfigi}[2][1]{\qf@check{#2}\qf@img[valign=c]{#1}{\linewidth}{#2}}
\makeatother
\newcommand{\papersec}[1]{\par\bigskip{\centering\Large\bfseries PAPER-#1\par}\medskip}
\newcommand{\partsec}[2]{\par\medskip\noindent{\large\bfseries PART-#1: #2}\par\smallskip}
\newcommand{\sect}[1]{\par\smallskip\noindent\textbf{#1}\par}
\newcommand{\qn}[1]{\par\smallskip\noindent\textbf{#1.}~}
\newcommand{\paraq}[2]{\par\smallskip\noindent\textbf{Paragraph for Questions #1 and #2}\par}
\newenvironment{qoptions}{\begin{itemize}[leftmargin=1.8em,itemsep=1pt,topsep=2pt,label={}]}{\end{itemize}}
\newcommand{\opt}[2]{\item[(#1)] #2}
% figure registry (generated by score2_tools.py finalize)
\qfigdef{p1_phy_s1_q02_1}{Score-I - Test 03 (07 Jan 2022)-img/p1_phy_s1_q02_1.png}{104.9}{0}
\qfigdef{p1_phy_s1_q03_1}{Score-I - Test 03 (07 Jan 2022)-img/p1_phy_s1_q03_1.png}{70.5}{0}
\qfigdef{p1_phy_s1_q05_1}{Score-I - Test 03 (07 Jan 2022)-img/p1_phy_s1_q05_1.png}{94.2}{0}
\qfigdef{p1_phy_s1_q08_1}{Score-I - Test 03 (07 Jan 2022)-img/p1_phy_s1_q08_1.png}{198.7}{0}
\qfigdef{p1_phy_s2_q09_1}{Score-I - Test 03 (07 Jan 2022)-img/p1_phy_s2_q09_1.png}{73.2}{0}
\qfigdef{p1_phy_s2_q10_1}{Score-I - Test 03 (07 Jan 2022)-img/p1_phy_s2_q10_1.png}{73.2}{0}
\qfigdef{p1_phy_s2_q11_1}{Score-I - Test 03 (07 Jan 2022)-img/p1_phy_s2_q11_1.png}{33.9}{0}
\qfigdef{p1_phy_s2_q11_2}{Score-I - Test 03 (07 Jan 2022)-img/p1_phy_s2_q11_2.png}{73.7}{0}
\qfigdef{p1_phy_s2_q11_3}{Score-I - Test 03 (07 Jan 2022)-img/p1_phy_s2_q11_3.png}{79.6}{0}
\qfigdef{p1_phy_s2_q11_4}{Score-I - Test 03 (07 Jan 2022)-img/p1_phy_s2_q11_4.png}{61.2}{0}
\qfigdef{p1_phy_s2_q12_1}{Score-I - Test 03 (07 Jan 2022)-img/p1_phy_s2_q12_1.png}{33.9}{0}
\qfigdef{p1_phy_s2_q12_2}{Score-I - Test 03 (07 Jan 2022)-img/p1_phy_s2_q12_2.png}{73.7}{0}
\qfigdef{p1_phy_s2_q12_3}{Score-I - Test 03 (07 Jan 2022)-img/p1_phy_s2_q12_3.png}{79.6}{0}
\qfigdef{p1_phy_s2_q12_4}{Score-I - Test 03 (07 Jan 2022)-img/p1_phy_s2_q12_4.png}{61.2}{0}
\qfigdef{p1_phy_s3_q02_1}{Score-I - Test 03 (07 Jan 2022)-img/p1_phy_s3_q02_1.png}{94.4}{0}
\qfigdef{p1_phy_s3_q03_1}{Score-I - Test 03 (07 Jan 2022)-img/p1_phy_s3_q03_1.png}{155.0}{0}
\qfigdef{p1_phy_s3_q06_1}{Score-I - Test 03 (07 Jan 2022)-img/p1_phy_s3_q06_1.png}{164.4}{0}
\qfigdef{p1_chem_s1_q04_1}{Score-I - Test 03 (07 Jan 2022)-img/p1_chem_s1_q04_1.png}{160.7}{0}
\qfigdef{p1_chem_s1_q06_1}{Score-I - Test 03 (07 Jan 2022)-img/p1_chem_s1_q06_1.png}{85.0}{0}
\qfigdef{p1_chem_s2_q09_2}{Score-I - Test 03 (07 Jan 2022)-img/p1_chem_s2_q09_2.png}{58.2}{0}
\qfigdef{p1_chem_s2_q09_3}{Score-I - Test 03 (07 Jan 2022)-img/p1_chem_s2_q09_3.png}{49.4}{0}
\qfigdef{p1_chem_s2_q10_1}{Score-I - Test 03 (07 Jan 2022)-img/p1_chem_s2_q10_1.png}{67.7}{0}
\qfigdef{p1_chem_s2_q10_2}{Score-I - Test 03 (07 Jan 2022)-img/p1_chem_s2_q10_2.png}{57.1}{0}
\qfigdef{p1_chem_s3_q04_1}{Score-I - Test 03 (07 Jan 2022)-img/p1_chem_s3_q04_1.png}{245.2}{0}
\qfigdef{p1_chem_s3_q05_1}{Score-I - Test 03 (07 Jan 2022)-img/p1_chem_s3_q05_1.png}{179.5}{0}
\qfigdef{p2_phy_s1_q01_2}{Score-I - Test 03 (07 Jan 2022)-img/p2_phy_s1_q01_2.png}{130.3}{0}
\qfigdef{p2_phy_s1_q02_1}{Score-I - Test 03 (07 Jan 2022)-img/p2_phy_s1_q02_1.png}{71.5}{0}
\qfigdef{p2_phy_s1_q04_1}{Score-I - Test 03 (07 Jan 2022)-img/p2_phy_s1_q04_1.png}{81.5}{0}
\qfigdef{p2_phy_s1_q05_1}{Score-I - Test 03 (07 Jan 2022)-img/p2_phy_s1_q05_1.png}{152.3}{0}
\qfigdef{p2_phy_s1_q06_1}{Score-I - Test 03 (07 Jan 2022)-img/p2_phy_s1_q06_1.png}{135.9}{0}
\qfigdef{p2_phy_s1_q07_1}{Score-I - Test 03 (07 Jan 2022)-img/p2_phy_s1_q07_1.png}{132.2}{0}
\qfigdef{p2_phy_s1_q09_1}{Score-I - Test 03 (07 Jan 2022)-img/p2_phy_s1_q09_1.png}{80.6}{0}
\qfigdef{p2_phy_s1_q11_A}{Score-I - Test 03 (07 Jan 2022)-img/p2_phy_s1_q11_A.png}{94.6}{0}
\qfigdef{p2_phy_s1_q13_1}{Score-I - Test 03 (07 Jan 2022)-img/p2_phy_s1_q13_1.png}{120.7}{0}
\qfigdef{p2_phy_s1_q14_1}{Score-I - Test 03 (07 Jan 2022)-img/p2_phy_s1_q14_1.png}{152.3}{0}
\qfigdef{p2_phy_s1_q15_1}{Score-I - Test 03 (07 Jan 2022)-img/p2_phy_s1_q15_1.png}{57.3}{0}
\qfigdef{p2_phy_s1_q16_1}{Score-I - Test 03 (07 Jan 2022)-img/p2_phy_s1_q16_1.png}{82.5}{0}
\qfigdef{p2_phy_s1_q17_1}{Score-I - Test 03 (07 Jan 2022)-img/p2_phy_s1_q17_1.png}{152.6}{0}
\qfigdef{p2_phy_s1_q18_1}{Score-I - Test 03 (07 Jan 2022)-img/p2_phy_s1_q18_1.png}{102.5}{0}
\qfigdef{p2_phy_s1_q19_1}{Score-I - Test 03 (07 Jan 2022)-img/p2_phy_s1_q19_1.png}{110.2}{0}
\qfigdef{p2_phy_s1_q20_1}{Score-I - Test 03 (07 Jan 2022)-img/p2_phy_s1_q20_1.png}{39.7}{0}
\qfigdef{p2_phy_s2_q01_1}{Score-I - Test 03 (07 Jan 2022)-img/p2_phy_s2_q01_1.png}{130.0}{0}
\qfigdef{p2_phy_s2_q02_1}{Score-I - Test 03 (07 Jan 2022)-img/p2_phy_s2_q02_1.png}{75.1}{0}
\qfigdef{p2_phy_s2_q04_1}{Score-I - Test 03 (07 Jan 2022)-img/p2_phy_s2_q04_1.png}{118.5}{0}
\qfigdef{p2_phy_s2_q05_1}{Score-I - Test 03 (07 Jan 2022)-img/p2_phy_s2_q05_1.png}{74.2}{0}
\qfigdef{p2_phy_s2_q06_1}{Score-I - Test 03 (07 Jan 2022)-img/p2_phy_s2_q06_1.png}{126.8}{0}
\qfigdef{p2_phy_s2_q07_1}{Score-I - Test 03 (07 Jan 2022)-img/p2_phy_s2_q07_1.png}{139.9}{0}
\qfigdef{p2_phy_s2_q08_1}{Score-I - Test 03 (07 Jan 2022)-img/p2_phy_s2_q08_1.png}{85.5}{0}
\qfigdef{p2_phy_s2_q09_1}{Score-I - Test 03 (07 Jan 2022)-img/p2_phy_s2_q09_1.png}{129.3}{0}
\qfigdef{p2_phy_s2_q10_1}{Score-I - Test 03 (07 Jan 2022)-img/p2_phy_s2_q10_1.png}{62.9}{0}
\qfigdef{p2_chem_s1_q08_1}{Score-I - Test 03 (07 Jan 2022)-img/p2_chem_s1_q08_1.png}{103.9}{0}
\qfigdef{p2_chem_s1_q10_1}{Score-I - Test 03 (07 Jan 2022)-img/p2_chem_s1_q10_1.png}{102.7}{0}
\qfigdef{p2_chem_s1_q15_1}{Score-I - Test 03 (07 Jan 2022)-img/p2_chem_s1_q15_1.png}{88.9}{0}
\qfigdef{p2_chem_s1_q15_2}{Score-I - Test 03 (07 Jan 2022)-img/p2_chem_s1_q15_2.png}{138.1}{0}
\qfigdef{p2_chem_s1_q15_3}{Score-I - Test 03 (07 Jan 2022)-img/p2_chem_s1_q15_3.png}{153.7}{0}
\qfigdef{p2_chem_s1_q16_1}{Score-I - Test 03 (07 Jan 2022)-img/p2_chem_s1_q16_1.png}{150.6}{0}
\qfigdef{p2_chem_s1_q16_2}{Score-I - Test 03 (07 Jan 2022)-img/p2_chem_s1_q16_2.png}{151.8}{0}
\qfigdef{p2_chem_s1_q17_1}{Score-I - Test 03 (07 Jan 2022)-img/p2_chem_s1_q17_1.png}{135.9}{0}
\qfigdef{p2_chem_s1_q18_1}{Score-I - Test 03 (07 Jan 2022)-img/p2_chem_s1_q18_1.png}{104.9}{0}
\qfigdef{p2_chem_s1_q18_A}{Score-I - Test 03 (07 Jan 2022)-img/p2_chem_s1_q18_A.png}{63.2}{0}
\qfigdef{p2_chem_s1_q18_B}{Score-I - Test 03 (07 Jan 2022)-img/p2_chem_s1_q18_B.png}{60.3}{0}
\qfigdef{p2_chem_s1_q18_C}{Score-I - Test 03 (07 Jan 2022)-img/p2_chem_s1_q18_C.png}{68.0}{0}
\qfigdef{p2_chem_s1_q18_D}{Score-I - Test 03 (07 Jan 2022)-img/p2_chem_s1_q18_D.png}{72.8}{0}
\qfigdef{p2_chem_s1_q19_1}{Score-I - Test 03 (07 Jan 2022)-img/p2_chem_s1_q19_1.png}{138.9}{0}
\qfigdef{p2_chem_s1_q19_A}{Score-I - Test 03 (07 Jan 2022)-img/p2_chem_s1_q19_A.png}{60.4}{0}
\qfigdef{p2_chem_s1_q19_B}{Score-I - Test 03 (07 Jan 2022)-img/p2_chem_s1_q19_B.png}{59.5}{0}
\qfigdef{p2_chem_s1_q19_C}{Score-I - Test 03 (07 Jan 2022)-img/p2_chem_s1_q19_C.png}{60.1}{0}
\qfigdef{p2_chem_s1_q19_D}{Score-I - Test 03 (07 Jan 2022)-img/p2_chem_s1_q19_D.png}{60.0}{0}
\qfigdef{p2_chem_s1_q20_1}{Score-I - Test 03 (07 Jan 2022)-img/p2_chem_s1_q20_1.png}{132.0}{0}
\qfigdef{p2_chem_s1_q20_A}{Score-I - Test 03 (07 Jan 2022)-img/p2_chem_s1_q20_A.png}{55.9}{0}
\qfigdef{p2_chem_s1_q20_B}{Score-I - Test 03 (07 Jan 2022)-img/p2_chem_s1_q20_B.png}{57.3}{0}
\qfigdef{p2_chem_s1_q20_C}{Score-I - Test 03 (07 Jan 2022)-img/p2_chem_s1_q20_C.png}{58.0}{0}
\qfigdef{p2_chem_s1_q20_D}{Score-I - Test 03 (07 Jan 2022)-img/p2_chem_s1_q20_D.png}{59.6}{0}
\qfigdef{p2_chem_s2_q01_1}{Score-I - Test 03 (07 Jan 2022)-img/p2_chem_s2_q01_1.png}{151.3}{0}
\qfigdef{p2_chem_s2_q02_1}{Score-I - Test 03 (07 Jan 2022)-img/p2_chem_s2_q02_1.png}{150.9}{0}
\qfigdef{p2_chem_s2_q03_1}{Score-I - Test 03 (07 Jan 2022)-img/p2_chem_s2_q03_1.png}{128.0}{0}
\qfigdef{p2_chem_s2_q04_1}{Score-I - Test 03 (07 Jan 2022)-img/p2_chem_s2_q04_1.png}{135.2}{0}
\qfigdef{p2_chem_s2_q07_1}{Score-I - Test 03 (07 Jan 2022)-img/p2_chem_s2_q07_1.png}{38.7}{0}
\qfigdef{p2_chem_s2_q07_2}{Score-I - Test 03 (07 Jan 2022)-img/p2_chem_s2_q07_2.png}{31.8}{0}

\begin{document}
\begin{center}{\Large\bfseries JEE (Main+Advanced) Score-I --- Test 03 (07 Jan 2022)}\end{center}
\begin{multicols}{2}
\papersec{1}
\partsec{1}{PHYSICS}
\sect{Section-I(i)}
\qn{1} A thin disc of radius $R$ carries charge $q$ uniformly distributed on its surface. Now disc is transformed into a solid sphere of radius $R$ which carries same amount of charge $q$ uniformly distributed through its volume. Choose the CORRECT statement(s).
\begin{qoptions}
\opt{A}{Work done by electrostatic force is equal to $\dfrac{q^2}{R\epsilon_0}\left[\dfrac{2}{3\pi^2}-\dfrac{3}{20\pi}\right]$}
\opt{B}{Ratio of self potential energy of disc to that of solid sphere is $\dfrac{40}{9\pi}$}
\opt{C}{Ratio of potential at centre of disc to potential at centre of solid sphere is $\dfrac{2}{\pi}$}
\opt{D}{Ratio of potential at centre of disc to potential at centre of solid sphere is $\dfrac{4}{3}$}
\end{qoptions}
\qn{2} A small dipole of dipole moment $p$ is placed at the origin oriented along the $y$-axis. Position vector $\vec r$ of point A makes angle $\theta$ with $z$ axis while projection of $\vec r$ on $xy$ plane makes angle $\phi$ with $x$ axis. Choose the CORRECT statement(s).
\qfig{p1_phy_s1_q02_1}
\begin{qoptions}
\opt{A}{Magnitude of electric field at point A is $\dfrac{P}{4\pi\epsilon_0 r^3}\sqrt{1+3\sin^2\theta\sin^2\phi}$}
\opt{B}{Electric potential at point A is $\dfrac{P\sin\theta\sin\phi}{4\pi\epsilon_0 r^2}$}
\opt{C}{Magnitude of electric field at point A is $\dfrac{P}{4\pi\epsilon_0 r^3}\sqrt{1+3\cos^2\theta\sin^2\phi}$}
\opt{D}{Electric potential at point A is $\dfrac{P\sin\theta\cos\phi}{4\pi\epsilon_0 r^2}$}
\end{qoptions}
\qn{3} A rough vertical circular ring of radius $R$, carrying a bead of mass $m$, rotates in its own plane about its centre with uniform angular velocity $\omega$. The coefficient of friction between ring and bead is $\mu$. Choose the correct statement(s).
\qfig{p1_phy_s1_q03_1}
\begin{qoptions}
\opt{A}{Minimum angular speed of ring such that bead does not slip on ring is $\sqrt{\dfrac{g}{\mu R}}\sqrt{1+\mu^2}$}
\opt{B}{Minimum angular speed of ring such that bead does not slip on ring is $\sqrt{\dfrac{g}{\mu R}}$}
\opt{C}{If bead does not slip on ring, maximum frictional force experienced by it is equal to $mg$}
\opt{D}{If bead does not slip on ring, minimum normal reaction force experienced by it is equal to $\left(|m\omega^2R-mg|\right)$}
\end{qoptions}
\qn{4} In Yuri Gagarin's first manned space flight in 1961, orbit was a closed orbit around Earth. The perigee and apogee were 181 km and 327 km above the Earth. Radius of Earth is 6400 km and $g=10\ \mathrm{ms^{-2}}$ at earth surface. Choose the CORRECT statement(s).
\begin{qoptions}
\opt{A}{Time period of orbit was around 5326 seconds}
\opt{B}{Time period of orbit was around 3400 seconds}
\opt{C}{Maximum speed in orbit was around $\dfrac{8\,\mathrm{km}}{\mathrm{sec}}$}
\opt{D}{Semi-major axis of orbit was 6654 km}
\end{qoptions}
\qn{5} Poor Dinanath (D) is thrown into a circular arena of radius $a$ containing a lion (L). Initially the lion is at the centre O of the arena while Dinanath is at $(a,0)$. Dinanath decides to run at his maximum speed $u$ around the perimeter. The lion responds by running at its maximum speed $v$ in such a way that it remains on the (moving) radius OD. Choose the correct statement(s).
\qfig{p1_phy_s1_q05_1}
\begin{qoptions}
\opt{A}{During the chase, distance between lion and Dinanath at time $t$ is $\left[a-\dfrac{va}{u}\sin\left(\dfrac{ut}{a}\right)\right]$}
\opt{B}{If $a=10$ m, $u=5$ m/s and $v=10$ m/s then lion catches Dinanath after $\dfrac{\pi}{3}$ s}
\opt{C}{Path followed by lion is circle}
\opt{D}{Path followed by lion is hyperbola}
\end{qoptions}
\qn{6} A square platform lies in horizontal X-Z plane such that its vertices are at $(8\,\mathrm{m}, 0\,\mathrm{m}, 16\,\mathrm{m})$, $(16\,\mathrm{m}, 0\,\mathrm{m}, 16\,\mathrm{m})$, $(8\,\mathrm{m}, 0\,\mathrm{m}, 24\,\mathrm{m})$ \& $(16\,\mathrm{m}, 0\,\mathrm{m}, 24\,\mathrm{m})$. A particle is projected with velocity $\left(v_2\hat i+25\hat j\right)$ m/s from origin w.r.t wind which blows at velocity $\left(v_1\hat k\right)$ m/s w.r.t. ground. At the same instant platform starts moving from rest with acceleration $\left(2\hat i+2.5\hat j\right)$ m s$^{-2}$. Acceleration due to gravity is $\left(-10\hat j\right)$ ms$^{-2}$. Choose the correct statement(s).
\begin{qoptions}
\opt{A}{Possible value of $v_2$ if particle hits the platform is $6\leqslant v_2\leqslant10$}
\opt{B}{Possible value of $v_1$ if particle hits the platform is $4\leqslant v_1\leqslant6$}
\opt{C}{Possible time of flight of particle is 6 sec.}
\opt{D}{Possible time of flight of particle is 4 sec.}
\end{qoptions}
\qn{7} A bead of mass $m$ can slide on a smooth circular wire of radius $a$, which is fixed in a vertical plane. The bead is connected to the highest point of the wire by a light spring of natural length $3a/2$ and spring constant $k$. P is lowest point on ring. Choose the CORRECT statement(s).
\begin{qoptions}
\opt{A}{For bead, P is stable equilibrium position if $k=\dfrac{2mg}{a}$}
\opt{B}{For bead, P is unstable equilibrium position if $k=\dfrac{2mg}{a}$}
\opt{C}{For bead, P is stable equilibrium position if $k=\dfrac{5mg}{a}$}
\opt{D}{For bead, P is unstable equilibrium position if $k=\dfrac{5mg}{a}$}
\end{qoptions}
\qn{8} The block has mass $M$ and rests on a surface for which the coefficients of static and kinetic friction are $\mu_s$ and $\mu_k$ respectively. A force $F=kt^2$ is applied to the cable. Given : $M=150$ kg, $k=60$ N/s$^2$, $\mu_s=0.5$, $\mu_k=0.4$. Which of the following is/are correct?
\qfig{p1_phy_s1_q08_1}
\begin{qoptions}
\opt{A}{The time at which block begins to slide is 2.5 sec}
\opt{B}{The time at which block begins to slide is 1.5 sec}
\opt{C}{Acceleration of block at $t=5$ sec is 16 m/s$^2$}
\opt{D}{Acceleration of block at $t=5$ sec is 6 m/s$^2$}
\end{qoptions}
\sect{Section-I(ii)}
\qn{9} Answer the following by appropriately matching the lists based on the information given in the paragraph

A rough track connects point O $(0\,\mathrm{m},0\,\mathrm{m})$ to point A$(1\,\mathrm{m},1\,\mathrm{m})$. External force $F$ always tangential to path pulls object of mass 2 kg from point O$(0\,\mathrm{m},0\,\mathrm{m})$ to point A$(1\,\mathrm{m},1\,\mathrm{m})$ very slowly. Coefficient of friction between block and track is 0.5
\qfig{p1_phy_s2_q09_1}
In List-I some physical quantity regarding motion of block is mentioned while List-II has its corresponding numerical value.
\begin{center}\small\setlength{\tabcolsep}{3pt}
\begin{tabular}{|c|p{3.4cm}|c|c|}
\hline
\multicolumn{2}{|c|}{List-I} & \multicolumn{2}{c|}{List-II}\\ \hline
(I) & Work done by frictional force on block magnitude wise in Joules & (P) & $10\sqrt2$\\ \hline
(II) & Normal reaction force between block and track at $x=\dfrac12$ m in Newton & (Q) & $20\sqrt{\dfrac23}$\\ \hline
(III) & Work done by external force $F$ on block magnitude wise in Joules & (R) & $10$\\ \hline
(IV) & Net reaction force between block and track at $x=\dfrac34$ m in Newton & (S) & $20\sqrt{\dfrac{5}{13}}$\\ \hline
\multicolumn{2}{|c|}{} & (T) & $5\sqrt{15}$\\ \hline
\multicolumn{2}{|c|}{} & (U) & $30$\\ \hline
\end{tabular}
\end{center}
If equation of track is $y=x^2$ then choose the correct combination from List I \& List II.
\begin{qoptions}
\opt{A}{I $\to$ S;II $\to$ P;III $\to$ Q;IV $\to$ T}
\opt{B}{I $\to$ U;II $\to$ P;III $\to$ T;IV $\to$ S}
\opt{C}{I $\to$ Q;II $\to$ U;III $\to$ S;IV $\to$ T}
\opt{D}{I $\to$ R;II $\to$ P;III $\to$ U;IV $\to$ S}
\end{qoptions}
\qn{10} Answer the following by appropriately matching the lists based on the information given in the paragraph

A rough track connects point O $(0\,\mathrm{m},0\,\mathrm{m})$ to point A$(1\,\mathrm{m},1\,\mathrm{m})$. External force $F$ always tangential to path pulls object of mass 2 kg from point O$(0\,\mathrm{m},0\,\mathrm{m})$ to point A$(1\,\mathrm{m},1\,\mathrm{m})$ very slowly. Coefficient of friction between block and track is 0.5
\qfig{p1_phy_s2_q10_1}
In List-I some physical quantity regarding motion of block is mentioned while List-II has its corresponding numerical value.
\begin{center}\small\setlength{\tabcolsep}{3pt}
\begin{tabular}{|c|p{3.4cm}|c|c|}
\hline
\multicolumn{2}{|c|}{List-I} & \multicolumn{2}{c|}{List-II}\\ \hline
(I) & Work done by frictional force on block magnitude wise in Joules & (P) & $10\sqrt2$\\ \hline
(II) & Normal reaction force between block and track at $x=\dfrac12$ m in Newton & (Q) & $20\sqrt{\dfrac23}$\\ \hline
(III) & Work done by external force $F$ on block magnitude wise in Joules & (R) & $10$\\ \hline
(IV) & Net reaction force between block and track at $x=\dfrac34$ m in Newton & (S) & $20\sqrt{\dfrac{5}{13}}$\\ \hline
\multicolumn{2}{|c|}{} & (T) & $5\sqrt{15}$\\ \hline
\multicolumn{2}{|c|}{} & (U) & $30$\\ \hline
\end{tabular}
\end{center}
If equation of track is $y^2=x$ then choose the correct combination from List I \& List II.
\begin{qoptions}
\opt{A}{I $\to$ R;II $\to$ Q;III $\to$ U;IV $\to$ T}
\opt{B}{I $\to$ S;II $\to$ Q;III $\to$ T;IV $\to$ U}
\opt{C}{I $\to$ P;II $\to$ Q;III $\to$ R;IV $\to$ S}
\opt{D}{I $\to$ U;II $\to$ Q;III $\to$ T;IV $\to$ R}
\end{qoptions}
\qn{11} Answer the following by appropriately matching the lists based on the information given in the paragraph

List-I includes physical situations where a charged object is shown. Points A \& B are marked List-II gives value of constant $N$.
\begin{center}\small\setlength{\tabcolsep}{3pt}
\begin{tabular}{|c|p{3.6cm}|c|c|}
\hline
\multicolumn{2}{|c|}{List-I} & \multicolumn{2}{c|}{List-II}\\ \hline
(I) & \qfigi{p1_phy_s2_q11_1}\par Infinite thin line charge of charge per unit length $\lambda=\dfrac{Q}{R}$ & (P) & $0$\\ \hline
(II) & \qfigi{p1_phy_s2_q11_2}\par Uniformly charged thin hollow hemisphere of charge $Q$ and radius $R$. B is centre of hemisphere and A is a point in equatorial plane & (Q) & $\dfrac{1}{\pi^2}$\\ \hline
(III) & \qfigi{p1_phy_s2_q11_3}\par Ring of radius $R$ carries charge $Q$ uniformly. B is centre of ring and A is on its axis & (R) & $\dfrac{\ln 2}{2\pi}$\\ \hline
(IV) & \qfigi{p1_phy_s2_q11_4}\par A ring of radius $R$ centred at origin carries charge $+Q$ on its one half and $-Q$ on other as shown in figure. B is centre of ring and point A is at $(0,0,R)$. & (S) & $\dfrac{1}{2\pi}$\\ \hline
\multicolumn{2}{|c|}{} & (T) & $\dfrac{\sqrt2-1}{4\sqrt2\pi}$\\ \hline
\multicolumn{2}{|c|}{} & (U) & $\dfrac{1}{8\pi}$\\ \hline
\end{tabular}
\end{center}
Electric field at point B is $\dfrac{NQ}{\epsilon_0R^2}$. Pick body from List-I and value of $N$ from List-II \& choose the correct combination.
\begin{qoptions}
\opt{A}{I $\to$ U;II $\to$ S;III $\to$ Q;IV $\to$ P}
\opt{B}{I $\to$ Q;II $\to$ P;III $\to$ R;IV $\to$ S}
\opt{C}{I $\to$ S;II $\to$ U;III $\to$ P;IV $\to$ Q}
\opt{D}{I $\to$ R;II $\to$ S;III $\to$ P;IV $\to$ Q}
\end{qoptions}
\qn{12} Answer the following by appropriately matching the lists based on the information given in the paragraph

List-I includes physical situations where a charged object is shown. Points A \& B are marked List-II gives value of constant $N$.
\begin{center}\small\setlength{\tabcolsep}{3pt}
\begin{tabular}{|c|p{3.6cm}|c|c|}
\hline
\multicolumn{2}{|c|}{List-I} & \multicolumn{2}{c|}{List-II}\\ \hline
(I) & \qfigi{p1_phy_s2_q12_1}\par Infinite thin line charge of charge per unit length $\lambda=\dfrac{Q}{R}$ & (P) & $0$\\ \hline
(II) & \qfigi{p1_phy_s2_q12_2}\par Uniformly charged thin hollow hemisphere of charge $Q$ and radius $R$. B is centre of hemisphere and A is a point in equatorial plane & (Q) & $\dfrac{1}{\pi^2}$\\ \hline
(III) & \qfigi{p1_phy_s2_q12_3}\par Ring of radius $R$ carries charge $Q$ uniformly. B is centre of ring and A is on its axis & (R) & $\dfrac{\ln 2}{2\pi}$\\ \hline
(IV) & \qfigi{p1_phy_s2_q12_4}\par A ring of radius $R$ centred at origin carries charge $+Q$ on its one half and $-Q$ on other as shown in figure. B is centre of ring and point A is at $(0,0,R)$. & (S) & $\dfrac{1}{2\pi}$\\ \hline
\multicolumn{2}{|c|}{} & (T) & $\dfrac{\sqrt2-1}{4\sqrt2\pi}$\\ \hline
\multicolumn{2}{|c|}{} & (U) & $\dfrac{1}{8\pi}$\\ \hline
\end{tabular}
\end{center}
Magnitude of potential difference between point A and B is $\dfrac{NQ}{\epsilon_0R}$. Pick body from List-I and value of $N$ from List-II \& choose the correct combination.
\begin{qoptions}
\opt{A}{I $\to$ R;II $\to$ P;III $\to$ T;IV $\to$ P}
\opt{B}{I $\to$ P;II $\to$ U;III $\to$ R;IV $\to$ T}
\opt{C}{I $\to$ S;II $\to$ T;III $\to$ Q;IV $\to$ P}
\opt{D}{I $\to$ Q;II $\to$ P;III $\to$ T;IV $\to$ S}
\end{qoptions}
\sect{Section-II}
\qn{1} A raindrop falls vertically through stationary mist, collecting mass as it falls. The raindrop remains spherical and the rate of mass gain is proportional to its speed and the square of its radius. If the drop starts from rest with a negligible radius, then find its uniform acceleration in m/s$^2$. ($g=10\,\mathrm{ms^{-2}}$)
\qn{2} From a thin uniform circular plate, part shown in figure is cutout and centre of mass of this part is at $\vec{r}_{cm}=\dfrac{k\,a\left(\hat i+\hat j\right)}{\pi}$. Find $k$.
\qfig{p1_phy_s3_q02_1}
\qn{3} In the given figure curved surface of frustum has no charge on it while circular ends have uniform surface charge density $\sigma$. Point charge $q$ is placed on axis of frustum. Net force exerted by frustum on charge $q$ is $x\times10^{-6}$ N. Find $x$.
\qfig{p1_phy_s3_q03_1}
\qn{4} An infinite line charge 8.85 C/m is present on $z$ axis. A thin wire of length $2L$ with its midpoint at $(0,L,0)$ is parallel to $x$ axis and carries charge 17.7 C/m. Coulomb force between these two wires is $x\times10^{12}$ N. Find $x$
\qn{5} Uniform charge of $8.85\times10^{-10}$ C/m$^2$ is distributed on surface of hollow cone of radius 10 cm. Interaction potential energy of cone with a proton placed at vertex of cone is $n\times10^{-19}$. Find $n$.
\qn{6} Gravitational field produced by half ring of mass $M$ \& radius $R$ at point P$(0,0,R)$ is $\dfrac{GMk}{R^2}$. Find $k$
\qfig{p1_phy_s3_q06_1}
\partsec{2}{CHEMISTRY}
\sect{Section-I(i)}
\qn{1} Identify the correct statement(s) regarding boron containing product (X), (Y) from the given reaction :

Colemanite $+$ Na$_2$CO$_3$ $\to$ (X) $+$ CaCO$_3$ $+$ (Z)

(Z) $\xrightarrow{\mathrm{CO_2}}$ (X) $+$ Na$_2$CO$_3$

(X) $\xrightarrow{\Delta}$ (Y) $+$ H$_2$O $+$ (Z)
\begin{qoptions}
\opt{A}{On heating (Y) with copper sulphate in a non-luminous flame, cuprous metaborate is obtained.}
\opt{B}{Compound (X) is a tetra-nuclear crystal.}
\opt{C}{The aqueous solution of (X) acts as a buffer.}
\opt{D}{All boron atoms in (X) are tetrahedral in shape.}
\end{qoptions}
\qn{2} Identify the correct statements regarding oxo-acids of phosphorus :
\begin{qoptions}
\opt{A}{All H$_4$P$_2$O$_6$, H$_4$P$_2$O$_7$ and H$_4$P$_2$O$_8$ contain P--O--P linkage.}
\opt{B}{Among H$_3$PO$_2$, H$_3$PO$_3$ and H$_3$PO$_4$ order of reducing character is H$_3$PO$_2>$ H$_3$PO$_3>$ H$_3$PO$_4$}
\opt{C}{Heating of phosphorus acid (H$_3$PO$_3$) and hypo-phosphorus acid (H$_3$PO$_2$) follows disproportionation.}
\opt{D}{Oxo-acids of P contains p$\pi$--p$\pi$ overlap in P--O bonds.}
\end{qoptions}
\qn{3} Indentify the correct statement(s) among given options regarding group-16 commpounds -
\begin{qoptions}
\opt{A}{Reducing power of dioxides of group-16 decreases from SO$_2$ to PoO$_2$ while oxidising power increases from SO$_3$ to PoO$_3$.}
\opt{B}{Both rhombic and monoclinic forms of sulphur are soluble in CS$_2$.}
\opt{C}{The bleaching property of SO$_2$ is because of its oxidising nature and it is permanent in nature.}
\opt{D}{The bleaching property of SO$_2$ is because of it's reducing nature and temporary in nature.}
\end{qoptions}
\qn{4} The diagram given below shows the enthalpy change for the chemical reactions between A$_2$, B$_2$, AB and AB$_2$. (A$_2$ \& B$_2$ are elements)
\qfig{p1_chem_s1_q04_1}
Select correct statement(s)
\begin{qoptions}
\opt{A}{The standard heat of formation of AB$_{2(g)}$ is 90 kJ/mole}
\opt{B}{Oxidation of A$_2$(g) by B$_2$(g) into AB$_2$(g) in standard condition is endothermic process}
\opt{C}{The standard heat of formation of AB$_{(g)}$ is 100 kJ/mole}
\opt{D}{Oxidation of AB (g) by B$_2$(g) into AB$_2$(g) in standard condition is endothermic process}
\end{qoptions}
\qn{5} 10 moles of a liquid is 50\% converted into its vapour at normal boiling point 273$^\circ$C. If latent heat of vaporization of the liquid is 273 L-atm/mole, then which of the following statement(s) is/are correct for this process.
\begin{qoptions}
\opt{A}{Work done by the system is 224 L-atm approximately}
\opt{B}{$\Delta$H is equal to 1365 L-atm}
\opt{C}{Entropy of surrounding decreases by 2.5 L-atm/K}
\opt{D}{$\Delta$U is equal to 1589 L-atm}
\end{qoptions}
\qn{6} 2 moles of a diatomic nonreacting ideal gas at temperature $T_1$ (state 1) is taken in a closed container. If this gas changes its state from 1 to 2, 3, 4 \& 5 by doing four different kinds of reversible processes in absence of non P--V work as given below. Considering only isobaric, isochoric, isothermal \& adiabatic processes, select the correct statement(s) or relationship(s). ($R=8.3$ J/mole/K, $\log 2=0.3$, $\log 3=0.48$)
\qfig{p1_chem_s1_q06_1}
\begin{qoptions}
\opt{A}{Molar heat capacity of gas for process $1\to3$ is greater than molar heat capacity of gas for process $1\to4$.}
\opt{B}{$\Delta S_{1\to4}\simeq28.67$ J/K}
\opt{C}{$\Delta$H for the process $1\to3$ is $7RT_1$}
\opt{D}{Heat exchange (Q) in process : $Q_{1\to5}<Q_{1\to2}<Q_{1\to3}$}
\end{qoptions}
\qn{7} Which of the following statement(s) is/are correct ?
\begin{qoptions}
\opt{A}{Gabriel phthalimide reaction is used to synthesize primary aliphatic amines only.}
\opt{B}{Hoffmann bromamide reaction can be used to synthesize primary aliphatic amines but not primary aromatic amines.}
\opt{C}{In Hoffmann ammonolysis 1$^\circ$, 2$^\circ$ and 3$^\circ$ amines can be synthesized.}
\opt{D}{Aniline can be synthesized by reduction of nitrobenzene using Sn/HCl.}
\end{qoptions}
\qn{8} Consider a set of reactions

A $\xrightarrow{\text{Reduction}}$ B $\xrightarrow{\mathrm{CHCl_3/KOH}}$ C $\xrightarrow{\text{Catalytic reduction}}$ N-methyl aniline

Which of the following statement(s) is/are incorrect ?
\begin{qoptions}
\opt{A}{Compound `A' can be nitrobenzene}
\opt{B}{Conversion of B $\to$ C is known as Reimer Tiemann reaction}
\opt{C}{Compound `C' is a cyanide}
\opt{D}{Compound `C' can undergo Hofmann bromamide reaction}
\end{qoptions}
\sect{Section-I(ii)}
\qn{9} Consider the compounds given in List I and their reaction/properties given in List II.
\begin{center}\small\setlength{\tabcolsep}{3pt}
\begin{tabular}{|c|p{2.3cm}|c|p{3.2cm}|}
\hline
\multicolumn{2}{|c|}{\begin{tabular}[c]{@{}c@{}}List-I\\(compounds)\end{tabular}} & \multicolumn{2}{c|}{\begin{tabular}[c]{@{}c@{}}List-II\\(Reaction/properties)\end{tabular}}\\ \hline
(I) & HCOOH & (P) & Liberates H$_2$ when reacts with Na.\\ \hline
(II) & H$_3$C -- COOH & (Q) & Can undergo Hell-Volhard-Zelinsky reaction.\\ \hline
(III) & \qfigi{p1_chem_s2_q09_2} & (R) & Can change the colour of litmus from blue to red.\\ \hline
(IV) & \qfigi{p1_chem_s2_q09_3} & (S) & Forms ester upon reaction with C$_2$H$_5$OH/H$^+$\\ \hline
 & & (T) & Can form amide upon treatment with NH$_3$/$\Delta$.\\ \hline
 & & (U) & CO is formed upon reaction with PCl$_5$.\\ \hline
\end{tabular}
\end{center}
Which of the following set is correctly matched ?
\begin{qoptions}
\opt{A}{I $\to$ P,S,Q}
\opt{B}{II $\to$ Q,T,U}
\opt{C}{III $\to$ P,Q,R}
\opt{D}{IV $\to$ P,R}
\end{qoptions}
\qn{10} Consider the compounds given in List I and their reaction/properties given in List II.
\begin{center}\small\setlength{\tabcolsep}{3pt}
\begin{tabular}{|c|p{2.3cm}|c|p{3.2cm}|}
\hline
\multicolumn{2}{|c|}{\begin{tabular}[c]{@{}c@{}}List-I\\(compounds)\end{tabular}} & \multicolumn{2}{c|}{\begin{tabular}[c]{@{}c@{}}List-II\\(Reaction/properties)\end{tabular}}\\ \hline
(I) & HCOOH & (P) & Liberates H$_2$ when reacts with Na.\\ \hline
(II) & H$_3$C -- COOH & (Q) & Can undergo Hell-Volhard-Zelinsky reaction.\\ \hline
(III) & \qfigi{p1_chem_s2_q10_1} & (R) & Can change the colour of litmus from blue to red.\\ \hline
(IV) & \qfigi{p1_chem_s2_q10_2} & (S) & Forms ester upon reaction with C$_2$H$_5$OH/H$^+$\\ \hline
 & & (T) & Can form amide upon treatment with NH$_3$/$\Delta$.\\ \hline
 & & (U) & CO is formed upon reaction with PCl$_5$.\\ \hline
\end{tabular}
\end{center}
Which of the following set is incorrectly matched ?
\begin{qoptions}
\opt{A}{I $\to$ U,P}
\opt{B}{II $\to$ T,S,Q}
\opt{C}{III $\to$ Q,U}
\opt{D}{IV $\to$ P, R}
\end{qoptions}
\qn{11} Generally in group-18, very rare reactions are observed but fluorides of Xe show variety of reactions. Mainly these fluorides behave like a replacement of fluorine in redox reactions, to control the vigourisity of chemical reaction.
\begin{center}\small\setlength{\tabcolsep}{3pt}
\begin{tabular}{|c|p{2.5cm}|c|p{3.2cm}|}
\hline
\multicolumn{2}{|c|}{\begin{tabular}[c]{@{}c@{}}List-I\\Chemical Reaction\end{tabular}} & \multicolumn{2}{c|}{\parbox{3.6cm}{\centering List-II\\(Properties of products formed/properties of reaction)}}\\ \hline
(I) & XeF$_2$ + H$_2$O $\to$ & (P) & Disproportionation\\ \hline
(II) & XeF$_2$ + H$_2$ $\to$ & (Q) & One of the product or intermediate product is used in etching of glass\\ \hline
(III) & XeF$_4$ + H$_2$O $\to$ & (R) & At least one gaseous product is formed.\\ \hline
(IV) & XeF$_6$ + H$_2$O\par $\xrightarrow[\text{medium}]{\text{in Alkaline}}$ & (S) & One of the product can act as an oxidising agent\\ \hline
 & & (T) & One of the product(s) is paramagnetic\\ \hline
 & & (U) & Redox reaction\\ \hline
\end{tabular}
\end{center}
Select option in which all the correct set of properties are given regarding above reactions :
\begin{qoptions}
\opt{A}{II $\to$ Q,R,T,U}
\opt{B}{II $\to$ P,R,S,T,U}
\opt{C}{IV $\to$ P,Q,R,S,T,U}
\opt{D}{IV $\to$ P,Q,S,T}
\end{qoptions}
\qn{12} Generally in group-18, very rare reactions are observed but fluorides of Xe show variety of reactions. Mainly these fluorides behave like a replacement of fluorine in redox reactions, to control the vigourisity of chemical reaction.
\begin{center}\small\setlength{\tabcolsep}{3pt}
\begin{tabular}{|c|p{2.5cm}|c|p{3.2cm}|}
\hline
\multicolumn{2}{|c|}{\begin{tabular}[c]{@{}c@{}}List-I\\Chemical Reaction\end{tabular}} & \multicolumn{2}{c|}{\parbox{3.6cm}{\centering List-II\\(Properties of products formed/properties of reaction)}}\\ \hline
(I) & XeF$_2$ + H$_2$O $\to$ & (P) & Disproportionation\\ \hline
(II) & XeF$_2$ + H$_2$ $\to$ & (Q) & One of the product or intermediate product is used in etching of glass\\ \hline
(III) & XeF$_4$ + H$_2$O $\to$ & (R) & At least one gaseous product is formed.\\ \hline
(IV) & XeF$_6$ + H$_2$O\par $\xrightarrow[\text{medium}]{\text{in Alkaline}}$ & (S) & One of the product can act as an oxidising agent\\ \hline
 & & (T) & One of the product(s) is paramagnetic\\ \hline
 & & (U) & Redox reaction\\ \hline
\end{tabular}
\end{center}
Select the incorrect set of properties regarding given reactions :
\begin{qoptions}
\opt{A}{III $\to$ P,Q,R,S,T,U}
\opt{B}{IV $\to$ P,R,S,T,U}
\opt{C}{I $\to$ Q,R,S,T,U}
\opt{D}{II $\to$ P,Q,R,T}
\end{qoptions}
\sect{Section-II}
\qn{1} In a polytropic expansion process ($PV^x=$ constant), a diatomic ideal gas does a work $\dfrac{y}{2}$ Joule by absorbing `$y$' Joule heat from surrounding. The value of `$x$' is-
\qn{2} Equilibrium constant ($K$) at 27$^\circ$C for the decomposition of poisonous gas COCl$_2$, represented by

COCl$_2$(g) $\rightleftharpoons$ CO(g) $+$ Cl$_2$(g) is $\dfrac{1}{2}\times10^{-7}$ bar

Find magnitude of $\Delta_r G$ (kJ/mole) for the following reaction at 27$^\circ$C

CO(g) $+$ Cl$_2$(g) $\rightleftharpoons$ COCl$_2$(g)

if mixture of 1 mole COCl$_2$(g), 1 mole CO(g) and 2 mole Cl$_2$(g) is taken in closed container and maintained at 1 bar pressure and 27$^\circ$C. ($R=8.3$ J/mole/K)
\qn{3} Identify the total number of correct statement regarding halogens and their compounds :

(I) Cl$_2$ turns, moist starch iodide paper blue,.

(II) Reaction of excess chlorine with ammonia produces nitrogen gas.

(III) I$_2$O$_5$ is used for the quantitative estimation of carbon mono-oxide.

(IV) The reducing nature of hydrogen halides increases from HF to HI.

(V) ClO$_2$ is paramagnetic and can undergo dimerisation easily.

(VI) The acidic strength of oxo-acids of chlorine increases with increase in oxidation state of chlorine.

(VII) Chlorine is used as a bleaching agent in textile industry.

(VIII) Na$_2$S$_2$O$_3$ is used as a antichlore in textile industries
\qn{4}
\qfig{p1_chem_s3_q04_1}
If $x=$ Double bond equivalent of product C

$y=$ Number of hydrogen atoms in product B

Then $x/y=\ldots$.
\qn{5}
\qfig{p1_chem_s3_q05_1}
if $p=$ molecular mass of group \fbox{R}

$q=$ molecular mass of group \fbox{R$'$}

where \fbox{R} and \fbox{R$'$} are alkyl groups

then minimum value of $|p-q|$ is
\qn{6} SNG (substitute natural gas) is a gaseous mixture containing methane gas that can be used as fuel for various purposes. If following reaction is one of the method to produce SNG
$$4\mathrm{CO(g)}+8\mathrm{H_2(g)}\to3\mathrm{CH_4(g)}+\mathrm{CO_2(g)}+2\mathrm{H_2O(\ell)},$$
then find magnitude of $\Delta_rH^\circ$ (kJ/mole) for above reaction based on $\Delta_rH^\circ$ data of following reactions.

\begin{center}
\begin{tabular}{lr}
Reactions & $\Delta_rH^\circ$ (kJ/mole)\\
$\mathrm{CO(g)}+\frac12\mathrm{O_2(g)}\to\mathrm{CO_2(g)}$ & $-300$\\
$\mathrm{C(graphite)}+2\mathrm{H_2(g)}\to\mathrm{CH_4(g)}$ & $-75$\\
$\mathrm{H_2(g)}+\frac12\mathrm{O_2(g)}\to\mathrm{H_2O(\ell)}$ & $-300$\\
$\mathrm{C(graphite)}+\frac12\mathrm{O_2(g)}\to\mathrm{CO(g)}$ & $-100$\\
\end{tabular}
\end{center}
\partsec{3}{MATHEMATICS}
\sect{Section-I(i)}
\qn{1} An insurance company believes that people can be divided into two classes: those who are accident prone and those who are not. The company's statistics show that an accident-prone person will have an accident at some time within a fixed 1-year period with probability 0.4, whereas this probability decreases to 0.2 for a person who is not accident prone. If we assume that 30 percent of the population is accident prone, then the probability that a new policyholder will have an accident within a year of purchasing a policy is $\dfrac{p}{q}$, $p$ and $q$ are co-prime numbers then
\begin{qoptions}
\opt{A}{$p+q=63$}
\opt{B}{$p+q=126$}
\opt{C}{$q-p=37$}
\opt{D}{$q-p=74$}
\end{qoptions}
\qn{2} If $(1+2x+3x^2)^{10}=a_0+a_1x+a_2x^2+\ldots+a_{20}x^{20}$, then
\begin{qoptions}
\opt{A}{$a_1+a_2=230$}
\opt{B}{$|a_1-a_2|=190$}
\opt{C}{$a_0+a_2+a_4+\ldots+a_{20}=\dfrac{6^{10}+2^{10}}{2}$}
\opt{D}{$a_1+a_3+\ldots+a_{19}=\dfrac{6^{10}-2^{10}}{2}$}
\end{qoptions}
\qn{3} The equation of the tangents to the circle $x^2+y^2-6x-4y+5=0$ which is inclined at an angle of $45^\circ$ with the $x$-axis is
\begin{qoptions}
\opt{A}{$x-y+3=0$}
\opt{B}{$x-y-3=0$}
\opt{C}{$x-y-5=0$}
\opt{D}{$x-y+5=0$}
\end{qoptions}
\qn{4} Given the family of lines, $a(3x+4y+6)+b(x+y+2)=0$. The line $L$ of the family situated at the greatest distance from the point $(2,3)$ has equation $\ell x+my+8=0$, then
\begin{qoptions}
\opt{A}{$\ell^2+m^2=25$}
\opt{B}{distance of origin from $L$ is $\dfrac85$}
\opt{C}{$\ell^2+m^2=5$}
\opt{D}{distance of origin from $L$ is $\dfrac{2\sqrt2}{\sqrt5}$}
\end{qoptions}
\qn{5} Suppose that we toss 2 fair dice. Let $E_1$ denote the event that the sum of digits on the dice is 6 and $F$ denote the event that the first die shows 4. Then
\begin{qoptions}
\opt{A}{$P(E_1\cap F)=\dfrac{1}{36}$}
\opt{B}{$P(E_1)P(F)=\dfrac{5}{216}$}
\opt{C}{$E_1$ and $F$ are independent events}
\opt{D}{$E_1$ and $F$ are not independent events}
\end{qoptions}
\qn{6} Let $T_r$ denote the $r^{\text{th}}$ term in the expansion of $\left[2^x+(1/4^x)\right]^n$ ($x\in\mathbb R$, $n\in\mathbb N$). If the ratio $T_3:T_2=7:1$ and sum of the binomial coefficients of second and third terms is 36, then
\begin{qoptions}
\opt{A}{$n=7$}
\opt{B}{$n=8$}
\opt{C}{$x=\dfrac12$}
\opt{D}{$x=-\dfrac13$}
\end{qoptions}
\qn{7} Number of ways in which 13 identical apples can be distributed among 3 persons so that no two persons receive equal number of apples and each one can receive none one or more, is equal to-
\begin{qoptions}
\opt{A}{$^{16}C_2$}
\opt{B}{$84$}
\opt{C}{$^{15}C_2$}
\opt{D}{$85$}
\end{qoptions}
\qn{8} If two lines represented by the equation $x^3+bx^2y+cxy^2+y^3=0$ are at right angles then
\begin{qoptions}
\opt{A}{$b+c=-1$}
\opt{B}{$b+c=-2$}
\opt{C}{$b+c=-3$}
\opt{D}{$b+c=0$}
\end{qoptions}
\sect{Section-I(ii)}
\qn{9} Answer the following by appropriately matching the lists based on the information given in the paragraph

If $S=0$ is a circle in standard form and $L=0$ is a straight line in the plane of the circle, then, in general, the equation $S+\lambda L=0$ represents a circle. If $L=0$ cuts $S=0$ in two points P and Q, then $S+\lambda L=0$ represents circles passing through P and Q. In particular, if $L=0$ touches the circle $S=0$, then $S+\lambda L=0$ also touches $L=0$.
\begin{center}\small\setlength{\tabcolsep}{3pt}
\begin{tabular}{|c|p{3.7cm}|c|p{2.5cm}|}
\hline
\multicolumn{2}{|c|}{List-I} & \multicolumn{2}{c|}{List-II}\\ \hline
(I) & The equation of the circle whose centre lies on the line $x+y-11=0$ and which passes through the intersection of the circle $x^2+y^2-3x+2y-4=0$ and the line $2x+5y+2=0$ is & (P) & $2(x^2+y^2)-10x-5y+1=0$\\ \hline
(II) & The equation of the circle touching the line $2x+3y+1=0$ at $(1,-1)$ and cutting orthogonally the circle having the line segment joining $(0,3)$ and $(-2,-1)$ as diameter is & (Q) & $x^2+y^2-9x-13y-10=0$\\ \hline
(III) & The equation of the circle passing through the intersection of circles $x^2+y^2=4$ and $x^2+y^2-2x-4y+4=0$ and touching the line $x+2y-5=0$ is & (R) & $x^2+y^2-x-2y=0$\\ \hline
(IV) & Suppose $ax+by+c=0$ where $a,b,c$ are in A.P. be normal to a family of circles. The equation of the circle of the family which intersects the circle $x^2+y^2-4x-4y-1=0$ orthogonally is & (S) & $x^2+y^2-2x+4y-3=0$\\ \hline
 & & (T) & $x^2+y^2+x+2y=0$\\ \hline
 & & (U) & $x^2+y^2+x+2y+8=0$\\ \hline
\end{tabular}
\end{center}
Which of the following is the only CORRECT combination ?
\begin{qoptions}
\opt{A}{I $\to$ Q}
\opt{B}{I $\to$ P}
\opt{C}{II $\to$ Q}
\opt{D}{II $\to$ S}
\end{qoptions}
\qn{10} Answer the following by appropriately matching the lists based on the information given in the paragraph

If $S=0$ is a circle in standard form and $L=0$ is a straight line in the plane of the circle, then, in general, the equation $S+\lambda L=0$ represents a circle. If $L=0$ cuts $S=0$ in two points P and Q, then $S+\lambda L=0$ represents circles passing through P and Q. In particular, if $L=0$ touches the circle $S=0$, then $S+\lambda L=0$ also touches $L=0$.
\begin{center}\small\setlength{\tabcolsep}{3pt}
\begin{tabular}{|c|p{3.7cm}|c|p{2.5cm}|}
\hline
\multicolumn{2}{|c|}{List-I} & \multicolumn{2}{c|}{List-II}\\ \hline
(I) & The equation of the circle whose centre lies on the line $x+y-11=0$ and which passes through the intersection of the circle $x^2+y^2-3x+2y-4=0$ and the line $2x+5y+2=0$ is & (P) & $2(x^2+y^2)-10x-5y+1=0$\\ \hline
(II) & The equation of the circle touching the line $2x+3y+1=0$ at $(1,-1)$ and cutting orthogonally the circle having the line segment joining $(0,3)$ and $(-2,-1)$ as diameter is & (Q) & $x^2+y^2-9x-13y-10=0$\\ \hline
(III) & The equation of the circle passing through the intersection of circles $x^2+y^2=4$ and $x^2+y^2-2x-4y+4=0$ and touching the line $x+2y-5=0$ is & (R) & $x^2+y^2-x-2y=0$\\ \hline
(IV) & Suppose $ax+by+c=0$ where $a,b,c$ are in A.P. be normal to a family of circles. The equation of the circle of the family which intersects the circle $x^2+y^2-4x-4y-1=0$ orthogonally is & (S) & $x^2+y^2-2x+4y-3=0$\\ \hline
 & & (T) & $x^2+y^2+x+2y=0$\\ \hline
 & & (U) & $x^2+y^2+x+2y+8=0$\\ \hline
\end{tabular}
\end{center}
Which of the following is the only INCORRECT combination ?
\begin{qoptions}
\opt{A}{I $\to$ Q}
\opt{B}{II $\to$ P}
\opt{C}{III $\to$ R}
\opt{D}{III $\to$ S}
\end{qoptions}
\qn{11} Answer the following by appropriately matching the lists based on the information given in the paragraph

Consider all possible permutation of all the letters of the word ``ALCHEMIST'' (No letter being repeated)
\begin{center}\small\setlength{\tabcolsep}{3pt}
\begin{tabular}{|c|p{3.7cm}|c|p{2.3cm}|}
\hline
\multicolumn{2}{|c|}{List-I} & \multicolumn{2}{c|}{List-II}\\ \hline
(I) & The number of different permutations containing the word ``CHEMIST'' is & (P) & $4!\;6!$\\ \hline
(II) & The number of different permutations in which the first and last positions are occupied by vowels & (Q) & $3!$\\ \hline
(III) & The number of different permutations in which none of the letters A, E, I occurs in the last five positions, is & (R) & $^5C_2\;3!\;6!$\\ \hline
(IV) & The number of different permutations in which the letters A, E, I occur only in odd positions, is & (S) & $3!\;7!$\\ \hline
 & & (T) & $^5C_3\;3!\;6!$\\ \hline
 & & (U) & $1$\\ \hline
\end{tabular}
\end{center}
Which of the following is the only CORRECT combination ?
\begin{qoptions}
\opt{A}{I $\to$ Q}
\opt{B}{I $\to$ R}
\opt{C}{II $\to$ P}
\opt{D}{II $\to$ T}
\end{qoptions}
\qn{12} Answer the following by appropriately matching the lists based on the information given in the paragraph

Consider all possible permutation of all the letters of the word ``ALCHEMIST'' (No letter being repeated)
\begin{center}\small\setlength{\tabcolsep}{3pt}
\begin{tabular}{|c|p{3.7cm}|c|p{2.3cm}|}
\hline
\multicolumn{2}{|c|}{List-I} & \multicolumn{2}{c|}{List-II}\\ \hline
(I) & The number of different permutations containing the word ``CHEMIST'' is & (P) & $4!\;6!$\\ \hline
(II) & The number of different permutations in which the first and last positions are occupied by vowels & (Q) & $3!$\\ \hline
(III) & The number of different permutations in which none of the letters A, E, I occurs in the last five positions, is & (R) & $^5C_2\;3!\;6!$\\ \hline
(IV) & The number of different permutations in which the letters A, E, I occur only in odd positions, is & (S) & $3!\;7!$\\ \hline
 & & (T) & $^5C_3\;3!\;6!$\\ \hline
 & & (U) & $1$\\ \hline
\end{tabular}
\end{center}
Which of the following is the only INCORRECT combination ?
\begin{qoptions}
\opt{A}{I $\to$ Q}
\opt{B}{II $\to$ S}
\opt{C}{III $\to$ R,T}
\opt{D}{III $\to$ P}
\end{qoptions}
\sect{Section-II}
\qn{1} If the point $P(h,k)$ lies on the line $2x+3y=5$ such that $|PA-PB|$ is maximum where $A(2,3)$ and $B(1,2)$ then the value of $(hk)$ is equal to
\qn{2} A fair coin is tossed 5 times. The probability of getting at least three heads on consecutive tosses is
\qn{3} If $\dfrac{1}{1!\,10!}+\dfrac{1}{3!\,8!}+\dfrac{1}{5!\,6!}+\ldots.+\dfrac{1}{11!\,1!}=\dfrac{2^K}{n!}$ where $n$ and $k$ are coprime then $\dfrac{n}{k}$ is
\qn{4} Ten different coins are tossed simultaneously. The probability of getting at least 7 heads is
\qn{5} Radius of the circle which touches the $x$-axis at $(12,0)$ and cuts of chord of length 10 unit on the positive $y$--axis is
\qn{6} If the number of three-element subsets drawn from the set $S=\{1,2,3,\ldots,29,30\}$ are such that the sum of the three elements is a multiple of 3 is $P$ then the value of $\dfrac{P}{100}$ is
\papersec{2}
\partsec{1}{PHYSICS}
\sect{Section-I}
\qn{1} An infinite line charge is placed along $y$-axis and two point charges are placed on the $x$-axis as shown in the figure. Find $x$ for which field on the positive $x$-axis is zero $(x\gg\ell)$ and $\dfrac{q\ell}{\lambda}=2$ :
\qfig{p2_phy_s1_q01_2}
\begin{qoptions}
\opt{A}{$2$}
\opt{B}{$1$}
\opt{C}{$\dfrac12$}
\opt{D}{$4$}
\end{qoptions}
\qn{2} A point charge $Q$ is located just above the centre of the flat face of hemisphere as shown in figure. The electric flux through the flat face and curved face of hemisphere are respectively
\qfig{p2_phy_s1_q02_1}
\begin{qoptions}
\opt{A}{$\dfrac{Q}{\epsilon_0},\ \dfrac{-Q}{2\epsilon_0}$}
\opt{B}{$\dfrac{-Q}{2\epsilon_0},\ \dfrac{Q}{2\epsilon_0}$}
\opt{C}{$\dfrac{Q}{\epsilon_0},\ \dfrac{-Q}{\epsilon_0}$}
\opt{D}{$\dfrac{-Q}{\epsilon_0},\ \dfrac{Q}{\epsilon_0}$}
\end{qoptions}
\qn{3} Three charged conducting metal spheres of radius $R_1$, $R_2$ and $R_3$ are connected together by very thin wires and they are separated by large distances from each other. If $R_1<R_2<R_3$, at equilibrium, which of the following sets of relations involving the electric field strength $E$ generated by a sphere at a point just outside the surface, its potential $V$ and its charge $Q$ must hold between the spheres ?
\begin{qoptions}
\opt{A}{$V_1=V_2=V_3$, $E_1<E_2<E_3$, $Q_1<Q_2<Q_3$}
\opt{B}{$V_1=V_2=V_3$, $E_1>E_2>E_3$, $Q_1<Q_2<Q_3$}
\opt{C}{$V_1<V_2<V_3$, $E_1<E_2<E_3$, $Q_1=Q_2=Q_3$}
\opt{D}{$V_1>V_2>V_3$, $E_1<E_2<E_3$, $Q_1>Q_2>Q_3$}
\end{qoptions}
\qn{4} Two concentric conducting shells of radius $R$ and $2R$ carry same charge $Q$. The potential energy of system is:
\qfig{p2_phy_s1_q04_1}
\begin{qoptions}
\opt{A}{$\dfrac{kQ^2}{2R}$}
\opt{B}{$\dfrac34\dfrac{kQ^2}{R}$}
\opt{C}{$\dfrac54\dfrac{kQ^2}{R}$}
\opt{D}{$\dfrac{kQ^2}{R}$}
\end{qoptions}
\qn{5} A graph of the $X$ component of the electric field as a function of $x$ in a region of space is shown. The $Y$ and $Z$ components of the electric field are zero in this region. If the electric potential is 10V at the origin, then find the potential at $x=2.0$ m.
\qfig{p2_phy_s1_q05_1}
\begin{qoptions}
\opt{A}{$10$V}
\opt{B}{$40$V}
\opt{C}{$-10$V}
\opt{D}{$30$V}
\end{qoptions}
\qn{6} A point $q$ is placed outside an uncharged conducting sphere as shown in the figure. Then:
\qfig{p2_phy_s1_q06_1}
\begin{qoptions}
\opt{A}{The potential of the sphere is zero.}
\opt{B}{The potential of the sphere is $\dfrac{kq}{r}$}
\opt{C}{If the sphere is earthed charge appearing on its surface is $-q$}
\opt{D}{If the sphere is earthed charge appearing on its surface is $\dfrac{r}{R}(-q)$}
\end{qoptions}
\qn{7} Two point electric dipoles with dipole moments $P_1$ and $P_2$ are separated by a large distance $r$ with their dipole axes mutually perpendicular as shown. The electric force between the dipoles is (where $K=\dfrac{1}{4\pi\epsilon_0}$)
\qfig{p2_phy_s1_q07_1}
\begin{qoptions}
\opt{A}{$\dfrac{6KP_1P_2}{r^4}$}
\opt{B}{$\dfrac{3KP_1P_2}{r^4}$}
\opt{C}{$\dfrac{2KP_1P_2}{r^4}$}
\opt{D}{$\dfrac{KP_1P_2}{r^4}$}
\end{qoptions}
\qn{8} The height at which the acceleration due to gravity decreases by 36\% of its value on the surface of the earth is (The radius of the earth is $R$)
\begin{qoptions}
\opt{A}{$\dfrac{R}{6}$}
\opt{B}{$\dfrac{R}{4}$}
\opt{C}{$\dfrac{R}{2}$}
\opt{D}{$\dfrac23R$}
\end{qoptions}
\qn{9} A satellite of mass $m_0$ is moving in an elliptical orbit around a planet of mass $M$. Planet is fixed at one of the foci of elliptical orbit. The eccentricity of elliptical orbit is $e$. Points A and B are called perihelion and aphelion respectively. The ratio of radius of curvature of elliptical orbit at perihelion to aphelion is
\qfig{p2_phy_s1_q09_1}
\begin{qoptions}
\opt{A}{$e$}
\opt{B}{$\dfrac{1-e}{1+e}$}
\opt{C}{$\dfrac1e$}
\opt{D}{$1$}
\end{qoptions}
\qn{10} For a particle moving in space with velocity $\vec v$, which of the following is incorrect?
\begin{qoptions}
\opt{A}{$\left|\dfrac{d\vec v}{dt}\right|\geqslant\dfrac{d|\vec v|}{dt}$}
\opt{B}{$\displaystyle\int_0^t|\vec v|dt\geqslant\left|\int_0^t\vec v\,dt\right|$}
\opt{C}{$|\Delta\vec v|\leqslant\displaystyle\int_{t_1}^{t_2}|\vec a|dt$}
\opt{D}{$\dfrac{d|\vec v|}{dt}=\dfrac{d\vec v}{dt}\cdot\left(\dfrac{\vec v}{v}\right)$}
\end{qoptions}
\qn{11} In the given figure, the angle of inclination of the inclined plane is $30^\circ$. The horizontal velocity $V_0$ so that the particle hits the inclined plane perpendicularly is
\qfig{p2_phy_s1_q11_A}
\begin{qoptions}
\opt{A}{$v_0=\sqrt{\dfrac{2gH}{3}}$}
\opt{B}{$v_0=\sqrt{\dfrac{2gH}{7}}$}
\opt{C}{$v_0=\sqrt{\dfrac{gH}{5}}$}
\opt{D}{$v_0=\sqrt{\dfrac{gH}{7}}$}
\end{qoptions}
\qn{12} Two boys are standing at the points A and B on the ground, where $AB=a$. The boy at B starts running in a direction perpendicular to AB with velocity $v_1$. The boy at A starts running simultaneously with constant velocity $v$ and catches the other in time $t$ then $t$ is
\begin{qoptions}
\opt{A}{$\dfrac{a}{\sqrt{v^2+v_1^2}}$}
\opt{B}{$\dfrac{a}{v+v_1}$}
\opt{C}{$\dfrac{a}{v-v_1}$}
\opt{D}{$\dfrac{a}{\sqrt{v^2-v_1^2}}$}
\end{qoptions}
\qn{13} All surfaces shown in figure are smooth. System is released with the spring unstretched. In equilibrium, compression in the spring will be
\qfig{p2_phy_s1_q13_1}
\begin{qoptions}
\opt{A}{$\dfrac{2mg}{k}$}
\opt{B}{$\dfrac{(M+m)g}{\sqrt2k}$}
\opt{C}{$\dfrac{mg}{\sqrt2k}$}
\opt{D}{$\dfrac{mg}{k}$}
\end{qoptions}
\qn{14} On a plank of length 36 m and mass 3 kg, a small block of mass 1 kg is kept as shown. Friction coefficient between the block and the plank is 0.3 where as that between the plank and the ground is 0.6. At some instant an impulse is given to both of them such that they acquire a common initial velocity of 21 m/s as shown. After what time will the block get off the plank?
\qfig{p2_phy_s1_q14_1}
\begin{qoptions}
\opt{A}{$3\sqrt2$ sec}
\opt{B}{$4$sec}
\opt{C}{$5$sec}
\opt{D}{$4\sqrt2$ sec}
\end{qoptions}
\qn{15} Two beads A and B of equal mass `$m$' are connected by a light string. The beads are positioned to move on a smooth ring in a vertical plane as shown in the figure. The tension in the string just after release is
\qfig{p2_phy_s1_q15_1}
\begin{qoptions}
\opt{A}{$\sqrt2\,mg$}
\opt{B}{$\dfrac{mg}{2}$}
\opt{C}{$\dfrac{mg}{4}$}
\opt{D}{$\dfrac{mg}{\sqrt2}$}
\end{qoptions}
\qn{16} Power versus time graph for a given force as shown in figure. Then work done by the force upto time $t$ ($\leq t_0$).
\qfig{p2_phy_s1_q16_1}
\begin{qoptions}
\opt{A}{First increases then decreases}
\opt{B}{First decreases then increases.}
\opt{C}{Always increases.}
\opt{D}{Always decreases}
\end{qoptions}
\qn{17} An external agent moves the block of mass $m$ slowly from A to B, along a smooth hill such that every time he applies the force tangentially. Find the work done by agent in this interval.
\qfig{p2_phy_s1_q17_1}
\begin{qoptions}
\opt{A}{$\dfrac{m^2g^2H^2}{L}$}
\opt{B}{$\dfrac{mgH^2}{L}$}
\opt{C}{$mg(H+L)$}
\opt{D}{$mgH$}
\end{qoptions}
\qn{18} A bob is connected to an inextensible string of length 5m \& other end of string is fixed. Bob is given a velocity $u$ at lowest point \& string slacks at a height of 8m from lowest point. The value of $u$ is (take $g=10\,\mathrm{m/s^2}$)
\qfig{p2_phy_s1_q18_1}
\begin{qoptions}
\opt{A}{$\sqrt{95}\,\mathrm{m/s}$}
\opt{B}{$\sqrt{190}\,\mathrm{m/s}$}
\opt{C}{$\sqrt{200}\,\mathrm{m/s}$}
\opt{D}{$\sqrt{250}\,\mathrm{m/s}$}
\end{qoptions}
\qn{19} Two smooth identical stationary spheres are kept touching each other on a smooth horizontal floor as shown. A third identical sphere moving horizontally with a constant speed hits both stationary spheres symmetrically. If after collision the third sphere moves in same direction with one fourth of its initial speed, the coefficient of restitution will be
\qfig{p2_phy_s1_q19_1}
\begin{qoptions}
\opt{A}{$2/3$}
\opt{B}{$1/3$}
\opt{C}{$1/4$}
\opt{D}{$1/6$}
\end{qoptions}
\qn{20} An inextensible string connects two masses $m$ and $4m$ and the system rests on a smooth horizontal floor. An impulse $J$ is imparted to B as shown in figure. The tension in the string in the subsequent motion is
\qfig{p2_phy_s1_q20_1}
\begin{qoptions}
\opt{A}{$\dfrac{J^2}{m\ell}$}
\opt{B}{$\dfrac{J^2}{5m\ell}$}
\opt{C}{$\dfrac{4J^2}{5m\ell}$}
\opt{D}{$\dfrac{J^2}{20m\ell}$}
\end{qoptions}
\sect{Section-II}
\qn{1} A point charge $q=8.85\times10^{-9}$ C is placed at $(0,d,0)$ where $d=1$m. An insulating surface ABCA is on the $x$-$z$ plane, which is bounded by a quarter circle AB and two equal length straight line BC and CA as shown in the figure. Electric flux through the surface ABCA in Nm$^2$/C is (nearest integer) Take $\epsilon_0=8.85\times10^{-12}\ \dfrac{\mathrm{C^2}}{\mathrm{Nm^2}}$
\qfig{p2_phy_s2_q01_1}
\qn{2} A charge particle of charge $q$ \& mass $m$ is projected towards the center of uniformly charged ring of radius $R$ \& carrying a charge $Q$ from a distance of $\sqrt3R$ from the center along the axis, find the minimum value of velocity of projection (in m/s) so that particle reaches to $x=-\infty$ where $x$--axis is along the axis of ring \& center of ring is at origin. (given $Q=6\ \mu$c, $q=3\mu$c, $m=2$gm, $R=1$m)
\qfig{p2_phy_s2_q02_1}
\qn{3} A rocket of mass 40kg has 360 kg of fuel. The exhaust velocity of the fuel is 2.0 km/sec. Calculate minimum rate of consumption of fuel in kg/s so that the rocket may rise from the ground ($g=10$ m/sec$^2$).
\qn{4} A circular track is made in a wedge \& a particle is released from the top most point of circular track. The wedge is kept on a smooth horizontal surface \& initially it is kept in contact with a vertical surface. Find the maximum height (in m) raised by particle from the lowest point of circular track ($R=2.2$m \& mass of particle \& wedge is same)
\qfig{p2_phy_s2_q04_1}
\qn{5} Linear charge density on the ring is given by $\lambda=\lambda_0\cos\theta$. If dipole moment of this semi-circular ring is $n\lambda_0R^2$, find $n$
\qfig{p2_phy_s2_q05_1}
\qn{6} Find the speed (in m/s) of 10 kg block at the end of 3 second.
\qfig{p2_phy_s2_q06_1}
\qn{7} A rod AB is shown in figure. End A of the rod is fixed on the ground. Block is moving with velocity $\sqrt3$ m/s towards right. Find the speed (in m/s) of end B of the rod when rod makes an angle of $60^\circ$ with the ground.
\qfig{p2_phy_s2_q07_1}
\qn{8} A cone made of insulating material has a total charge $Q=3$mC spread uniformly over its sloping surface. Calculate the energy (in J) required to take a charge $q=5\ \mu$C from infinity to apex A (half angle of cone $=37^\circ$) (Radius $R=2$m)
\qfig{p2_phy_s2_q08_1}
\qn{9} A spaceship is sent to investigate a planet of mass $M$ and radius $R$. While hanging motionless in space at a distance $5R$ from the centre of the planet, the spaceship fires an instrument package with speed $v_0$ as shown in the figure. The package has mass $m$, which is much smaller than the mass of the spaceship. The angle $\theta$ for which the package just grazes the surface of the planet is $\sin^{-1}\left[\dfrac15\sqrt{1+\dfrac{kGM}{5v_0^2R}}\right]$. Find the value of $k$.
\qfig{p2_phy_s2_q09_1}
\qn{10} A ball A of mass $M=4$ kg is suspended by a vertical string. Another ball `B' of mass $m=1$ kg moving with a velocity $u=5.8$ m/s at an angle $\theta=53^\circ$ from vertical, collides elastically with the ball A as shown. Find the velocity of ball B, after collision in m/s.
\qfig{p2_phy_s2_q10_1}
\partsec{2}{CHEMISTRY}
\sect{Section-I}
\qn{1} Which of the following is incorrect regarding B$_3$N$_3$H$_6$ :-
\begin{qoptions}
\opt{A}{It is isoelectronic and isostructural with Benzene.}
\opt{B}{It is non - polar Compounds with non-polar bonds.}
\opt{C}{It shows electrophilic addition reaction}
\opt{D}{Borazine hydrolyzed at high temperature to give Ammonia gas, Boric acid and hydrogen gas}
\end{qoptions}
\qn{2} Which order is incorrect :-
\begin{qoptions}
\opt{A}{H$_2$O $>$ H$_2$Te $>$ H$_2$Se $>$ H$_2$S (Boiling point)}
\opt{B}{N$_2$O $<$ NO $<$ N$_2$O$_3$ $<$ N$_2$O$_4$ $<$ N$_2$O$_5$ (Acidic strength)}
\opt{C}{Br$_2>$ Cl$_2>$ F$_2$ (Boiling point)}
\opt{D}{Xe $<$ Kr $<$ Ar $<$ Ne $<$ He (Ease of liquification)}
\end{qoptions}
\qn{3} Oxyacid of phosphorous which contains P-P bond but no P - H bond is.
\begin{qoptions}
\opt{A}{Pyrophosphorous acid}
\opt{B}{Hypophosphoric acid}
\opt{C}{Pyrophosphoric acid}
\opt{D}{cyclotrimeta phosphoric acid}
\end{qoptions}
\qn{4} Which of the following statements are not correct ?
\begin{qoptions}
\opt{A}{O--O bond length is higher in O$_2$F$_2$ as compared to H$_2$O$_2$.}
\opt{B}{Three S - S linkages are present in a molecule of tetrathionic acid}
\opt{C}{Ozone turns starch-iodide paper Blue}
\opt{D}{Fluorapatite is a mineral of phosphorous}
\end{qoptions}
\qn{5} Which of the following statement is incorrect regarding allotropes ?
\begin{qoptions}
\opt{A}{Rhombic sulphur is the most stable at room temperature.}
\opt{B}{White phosphorus is the most reactive among allotropes of phosphorous.}
\opt{C}{Yellow phosphorus is soluble in CS$_2$ while red phosphorus is not}
\opt{D}{Graphite have sp$^2$ while diamond has sp hybridised `C' atom.}
\end{qoptions}
\qn{6} Incorrect statements regarding Xenon Fluorides --
\begin{qoptions}
\opt{A}{XeF$_4$ is square planar and acts as Fluoride donor with PF$_5$}
\opt{B}{XeF$_2$, XeF$_4$, and XeF$_6$ are colourless crystalline solids.}
\opt{C}{XeF$_6$ reacts with RbF to form Rb[XeF$_7$]}
\opt{D}{XeOF$_4$ and XeO$_2$F$_2$ are obtained by Complete hydrolysis of XeF$_6$}
\end{qoptions}
\qn{7} Select the redox reaction among given

(I) XeF$_6$ + NaF $\longrightarrow$

(II) XeF$_4$ + O$_2$F$_2$ $\longrightarrow$

(III) KI + H$_2$SO$_4$ (conc) $\xrightarrow{\Delta}$

(IV) P + H$_2$SO$_4$ (conc) $\xrightarrow{\Delta}$

(V) K$_4$[Fe(CN)$_6$] + H$_2$SO$_4$(dil.) $\xrightarrow{\Delta}$
\begin{qoptions}
\opt{A}{I , II, III}
\opt{B}{II, III, IV}
\opt{C}{I, II, III, IV}
\opt{D}{I, II, III, IV, V}
\end{qoptions}
\qn{8} An ideal gas is taken around the cycle ABCA as shown in P-V diagram. The value of $q$ for the cycle is equal to :
\qfig{p2_chem_s1_q08_1}
\begin{qoptions}
\opt{A}{$7\,P_1V_1$}
\opt{B}{$-5\,P_1V_1$}
\opt{C}{$5\,P_1V_1$}
\opt{D}{$-7\,P_1V_1$}
\end{qoptions}
\qn{9} Starting with the same initial conditions, one mole of an ideal monoatomic gas expands reversibly from volume $V_1$ to $V_2$ in three different ways. The work done by gas is $|\Delta W_1|$ if the process is isothermal, $|\Delta W_2|$ if isobaric and $|\Delta W_3|$ if adiabatic, then
\begin{qoptions}
\opt{A}{$|\Delta W_1|>|\Delta W_2|>|\Delta W_3|$}
\opt{B}{$|\Delta W_2|>|\Delta W_1|>|\Delta W_3|$}
\opt{C}{$|\Delta W_2|>|\Delta W_3|>|\Delta W_1|$}
\opt{D}{$|\Delta W_1|>|\Delta W_3|>|\Delta W_2|$}
\end{qoptions}
\qn{10} An ideal gas is taken from the same initial pressure $P_1$ to the same final pressure $P_2$ by three different processes. If it is known that point 1 corresponds to a reversible adiabatic and point 2 corresponds to a single stage adiabatic then
\qfig{p2_chem_s1_q10_1}
\begin{qoptions}
\opt{A}{Point 3 may be a two stage adiabatic.}
\opt{B}{The average K.E. of the gas is maximum at point 1}
\opt{C}{Work done by surrounding in reaching point number `3' will be minimum}
\opt{D}{Point 3 and point 4 may lie along an isotherm.}
\end{qoptions}
\qn{11} Two mole of an ideal gas is expanded irreversibly and isothermally at 37$^\circ$C until its volume is doubled and 3.41 kJ heat is absorbed from surrounding. $\Delta S_{\text{total}}$ (system $+$ surrounding) is :
\begin{qoptions}
\opt{A}{$-0.52$ J/K}
\opt{B}{$0.52$ J/K}
\opt{C}{$22.52$ J/K}
\opt{D}{$0$ J/K}
\end{qoptions}
\qn{12} In which of the following cases does entropy increase?
\begin{qoptions}
\opt{A}{C$_2$H$_4$(g) + H$_2$(g) $\longrightarrow$ C$_2$H$_6$(g)}
\opt{B}{N$_2$O$_4$(g) $\longrightarrow$ N$_2$O$_4$(s)}
\opt{C}{Graphite $\longrightarrow$ diamond}
\opt{D}{Combustion of glucose(s) at 110$^\circ$C}
\end{qoptions}
\qn{13} Steam undergoes decomposition at high temperature as per the reaction :

H$_2$O(g) $\rightleftharpoons$ H$_2$(g) + $\dfrac12$ O$_2$(g) , $\Delta$H$^\circ=200$ kJ mol$^{-1}$ $\Delta$S$^\circ=40$ kJ mol$^{-1}$.

The temperature at which equilibrium constant is unity is :
\begin{qoptions}
\opt{A}{3000 Kelvin}
\opt{B}{5000 Kelvin}
\opt{C}{5333 Kelvin}
\opt{D}{5 Kelvin}
\end{qoptions}
\qn{14} Calculate value of `$\ln K_{eq}$' for the reaction at 250 K.

N$_2$O$_4$(g) $\rightleftharpoons$ 2NO$_2$(g)

Given, $\Delta$H$^\circ_f$ (NO$_2$)g $=+40.407$ kJ / mol

$\Delta$H$^\circ_f$ (N$_2$O$_4$)g $=+70$ kJ / mol

$\Delta$S$^\circ$r $=10$ JK$^{-1}$
\begin{qoptions}
\opt{A}{4}
\opt{B}{$-4$}
\opt{C}{1.2}
\opt{D}{$-1.2$}
\end{qoptions}
\qn{15} Given two methods are suggested for the preparation of
\qfig{p2_chem_s1_q15_1}
Method-I :
\qfig{p2_chem_s1_q15_2}
Method-II :
\qfig{p2_chem_s1_q15_3}
Which is most appropriate method to carry out the conversion.
\begin{qoptions}
\opt{A}{Method I}
\opt{B}{Method II}
\opt{C}{Both method I and method II}
\opt{D}{Neither method I nor method II}
\end{qoptions}
\qn{16} P \& Q can be distinguished by :
\qfig{p2_chem_s1_q16_1}
\qfig{p2_chem_s1_q16_2}
\begin{qoptions}
\opt{A}{NaHSO$_3$}
\opt{B}{NaHCO$_3$}
\opt{C}{[Ag(NH$_3$)$_2$]$^+$}
\opt{D}{Br$_2$/CCl$_4$}
\end{qoptions}
\qn{17} P \& Q are major product of following reaction
\qfig{p2_chem_s1_q17_1}
Reagent which convert `Q' into `P' is / are :-
\begin{qoptions}
\opt{A}{Br$_2$ / NaOH}
\opt{B}{NaNO$_2$ / HCl}
\opt{C}{H$_3$O$^+$}
\opt{D}{None of these}
\end{qoptions}
\qn{18} What is the product of the following reaction?
\qfig{p2_chem_s1_q18_1}
\begin{qoptions}
\opt{A}{\qfigi{p2_chem_s1_q18_A}}
\opt{B}{\qfigi{p2_chem_s1_q18_B}}
\opt{C}{\qfigi{p2_chem_s1_q18_C}}
\opt{D}{\qfigi{p2_chem_s1_q18_D}}
\end{qoptions}
\qn{19}
\qfig{p2_chem_s1_q19_1}
Product A is :
\begin{qoptions}
\opt{A}{\qfigi{p2_chem_s1_q19_A}}
\opt{B}{\qfigi{p2_chem_s1_q19_B}}
\opt{C}{\qfigi{p2_chem_s1_q19_C}}
\opt{D}{\qfigi{p2_chem_s1_q19_D}}
\end{qoptions}
\qn{20} What is the product of the following synthesis?
\qfig{p2_chem_s1_q20_1}
\begin{qoptions}
\opt{A}{\qfigi{p2_chem_s1_q20_A}}
\opt{B}{\qfigi{p2_chem_s1_q20_B}}
\opt{C}{\qfigi{p2_chem_s1_q20_C}}
\opt{D}{\qfigi{p2_chem_s1_q20_D}}
\end{qoptions}
\sect{Section-II}
\qn{1} How many of the following give positive isocyanide test
\qfig{p2_chem_s2_q01_1}
\qn{2}
\qfig{p2_chem_s2_q02_1}
no of products including stereo isomerism
\qn{3}
\qfig{p2_chem_s2_q03_1}
How many product will be formed when above compound undergo de-carboxylation?
\qn{4} Total number of enol possible for the compound formed during given reaction will be (including stereoisomer):
\qfig{p2_chem_s2_q04_1}
\qn{5} How many total number of following orders are correct according to their given property ?

(I) Cl$_2>$ F$_2$ (Bond energy)

(II) HOCl $>$ HOBr (Acidic strength)

(III) HClO $<$ HClO$_2<$ HClO$_3<$ HClO$_4$ (thermal stability)

(IV) SF$_6<$ SeF$_6<$ TeF$_6$ (ease of hydrolysis)

(V) NH$_3>$ SbH$_3>$ AsH$_3>$ PH$_3$ (Melting point)

(VI) C $>$ Si $>$ Ge $>$ Sn $>$ Pb (IE$_1$)
\qn{6} Find the number of silicates in which only two corners per tetrahedron are shared -

Pyrosilicate, cyclic silicate, Double chain silicate, single chain silicate, 3D silicate, Sheet silicate.
\qn{7} Find the number of reagents which would lead to symmetrical cleavage of diborane -

NH$_3$ , CH$_3$NH$_2$, (CH$_3$)$_2$NH, \qfigi{p2_chem_s2_q07_1}, (CH$_3$)$_3$N, NaH, \qfigi{p2_chem_s2_q07_2}
\qn{8} How many of the following are correctly matched ?

(i) Work - state function.

(ii) Heat capacity - path function, extensive.

(iii) Specific internal energy - state function, intensive.

(iv) Molar Gibb's free energy - state function, extensive.

(v) Refractive index - extensive.

(vi) Dielectric constant - extensive.

(vii) Molar conductivity - intensive.

(viii) EMF of cell - intensive.
\qn{9} Certain amount of a gas confined in a piston-filled cylinder is heated from 27$^\circ$C to 127$^\circ$C keeping its pressure constant and the gas expanded against a constant pressure doing 4.157 kJ of work on surroundings. Find the number of moles of gas present in the cylinder ?
\qn{10} Standard enthalpy of formation of how many of the following is zero at 298 K.

Br$_2$(g), P(white), P(black), S(monoclinic), Cu(s), H$_2$S(g), H$^+$(aq), C(diamond),CS$_2$($\ell$)
\partsec{3}{MATHEMATICS}
\sect{Section-I}
\qn{1} If the point $P\left(\dfrac{3k}{2},1\right)$ lies between the two different lines $x+y=k$ and $x+y=2k$ then the least positive integral value of $k$ is
\begin{qoptions}
\opt{A}{$3$}
\opt{B}{$4$}
\opt{C}{$0$}
\opt{D}{$5$}
\end{qoptions}
\qn{2} In how many ways the allocation of three teams consisting of equal number of persons can be made for Swach Bharat Abhiyan using 18 persons to three team leaders Milkha, Amitabh \& Sachin -
\begin{qoptions}
\opt{A}{$\dfrac{18!}{(6!)^3\,3!}$}
\opt{B}{$\dfrac{18!}{(6!)^3}$}
\opt{C}{$\dfrac{18!}{(6!)^3}\cdot3!$}
\opt{D}{$\dfrac{18!}{(6!)^3\cdot(3!)^6}$}
\end{qoptions}
\qn{3} Given the line $x+y=3$ cuts the curve $x^2+y^2=a^2$ at two distinct points A and B such that OA is perpendicular to OB (where O is origin), then
\begin{qoptions}
\opt{A}{$a^2=3$}
\opt{B}{$a^2=9$}
\opt{C}{$a^2=81$}
\opt{D}{no such $a$ exist}
\end{qoptions}
\qn{4} Four fair dice $D_1$, $D_2$, $D_3$ and $D_4$, each having six faces numbered 1, 2, 3, 4, 5 and 6 are rolled simultaneously. The probability that $D_4$ shows a number appearing on one of $D_1$, $D_2$ and $D_3$ is -
\begin{qoptions}
\opt{A}{$\dfrac{91}{216}$}
\opt{B}{$\dfrac{108}{216}$}
\opt{C}{$\dfrac{125}{216}$}
\opt{D}{$\dfrac{127}{216}$}
\end{qoptions}
\qn{5} For all values of $\phi$, the lines represented by the equation
{\footnotesize\[\begin{gathered}(2\cos\phi+3\sin\phi)\,x+(3\cos\phi-5\sin\phi)\,y\\-(5\cos\phi-2\sin\phi)=0\end{gathered}\]}
\begin{qoptions}
\opt{A}{Pass through either of the two fixed points}
\opt{B}{Pass through the point $(4,5)$}
\opt{C}{pass through a fixed point whose reflection with respect to the line $x+y=\sqrt2$ is $\left(\sqrt2-1,\sqrt2-1\right)$}
\opt{D}{Pass through origin}
\end{qoptions}
\qn{6} Given 2 events A and B. If odd against A are as $2:1$ and odds in favour of $A\cup B$ are $3:1$ then range of $P(B)$ is
\begin{qoptions}
\opt{A}{$\left[\dfrac12,\dfrac5{12}\right]$}
\opt{B}{$\left[\dfrac12,\dfrac34\right]$}
\opt{C}{$\left[\dfrac5{12},\dfrac34\right]$}
\opt{D}{$[0,1]$}
\end{qoptions}
\qn{7} If a circle (passes through the point $(a,b)$) and cuts the circle $x^2+y^2=4$ orthogonally, then the equation of the locus of its centre is
\begin{qoptions}
\opt{A}{$2ax+2by-(a^2+b^2+4)=0$}
\opt{B}{$2ax+2by-(a^2-b^2+4)=0$}
\opt{C}{$x^2+y^2-3ax-4by+(a^2+b^2-4)=0$}
\opt{D}{$x^2+y^2-2ax-3by+(a^2-b^2-4)=0$}
\end{qoptions}
\qn{8} A man has 3 pairs of black socks and 2 pairs of brown socks, kept together in a box. If he dressed hurriedly in the dark, the probability that after he has put on a black sock, he will then put a brown sock is -
\begin{qoptions}
\opt{A}{$\dfrac13$}
\opt{B}{$\dfrac23$}
\opt{C}{$\dfrac4{15}$}
\opt{D}{$\dfrac{11}{15}$}
\end{qoptions}
\qn{9} The numerically greatest term in the expansion of $(2x+5y)^{34}$ when $x=3$ \& $y=2$ is
\begin{qoptions}
\opt{A}{$T_{21}$}
\opt{B}{$T_{22}$}
\opt{C}{$T_{23}$}
\opt{D}{$T_{24}$}
\end{qoptions}
\qn{10} $A$ is a set containing 7 elements. A subset $A_1$ of $A$ is chosen. Then set $A$ is reconstructed by replacing the elements of $A_1$. Now another subset $A_2$ of $A$ is chosen. In this way 4 Subsets $A_1$, $A_2$, $A_3$, $A_4$ are formed. The number of ways of choosing $A_1$, $A_2$, $A_3$, $A_4$ So that these subsets are pairwise disjoint, is
\begin{qoptions}
\opt{A}{$5^7$}
\opt{B}{$7^5$}
\opt{C}{$4^7$}
\opt{D}{$840$}
\end{qoptions}
\qn{11} The coefficient of $x^6$ in the expansion of $\left(x^2+\dfrac{4}{x^2}+3\right)^5$ must be -
\begin{qoptions}
\opt{A}{$90$}
\opt{B}{$110$}
\opt{C}{$100$}
\opt{D}{$120$}
\end{qoptions}
\qn{12} If three students A, B, C independently solve a problem with probabilities $\dfrac13$, $\dfrac14$ and $\dfrac15$ respectively, then the probability that the problem will be solved is
\begin{qoptions}
\opt{A}{$\dfrac35$}
\opt{B}{$\dfrac45$}
\opt{C}{$\dfrac25$}
\opt{D}{$\dfrac{47}{60}$}
\end{qoptions}
\qn{13} If $^{10}C_0+{}^{11}C_1+{}^{12}C_2+\cdots\cdots\cdots+{}^{100}C_{90}={}^xC_y$, then $(x,y)$ is equal to :-
\begin{qoptions}
\opt{A}{$(101,91)$}
\opt{B}{$(101,90)$}
\opt{C}{$(101,10)$}
\opt{D}{$(102,91)$}
\end{qoptions}
\qn{14} Let the algebraic sum of the perpendicular distances from $(3,0),(0,3)$ and $(2,2)$ to a variable line is zero, and the line passes through a fixed point whose co-ordinates are $(a,b)$. Then the value of $a-b$ is
\begin{qoptions}
\opt{A}{$0$}
\opt{B}{$1$}
\opt{C}{$2$}
\opt{D}{$\dfrac{10}{3}$}
\end{qoptions}
\qn{15} Let A,B and C be three events such that $P\left(A\cap\bar B\cap\bar C\right)=0.6$, $P(A)=0.8$ and $P\left(\bar A\cap B\cap C\right)=0.1$, then the value of $P$(atleast two among A,B and C) equals-
\begin{qoptions}
\opt{A}{$0.1$}
\opt{B}{$0.2$}
\opt{C}{$0.3$}
\opt{D}{$0.5$}
\end{qoptions}
\qn{16} The values of $\lambda$ for which the set
{\footnotesize\[\begin{gathered}S=\left\{(x,y):x^2+y^2+2x\leqslant1\right\}\\\cap\left\{(x,y):x-y+\lambda\geqslant0\right\}\end{gathered}\]}
contains only one point is
\begin{qoptions}
\opt{A}{$(-\infty,-1]\cup[3,\infty)$}
\opt{B}{$\{-1,3\}$}
\opt{C}{$\{-3\}$}
\opt{D}{$\{1\}$}
\end{qoptions}
\qn{17} Let a variable line intersects the curve $3x^2-y^2-2x+4y=0$ at A and B such that OA is perpendicular to OB (where O is origin). Then the variable line passes through the fixed point
\begin{qoptions}
\opt{A}{$(2,3)$}
\opt{B}{$(1,-2)$}
\opt{C}{$(3,4)$}
\opt{D}{$\left(\dfrac12,-\dfrac23\right)$}
\end{qoptions}
\qn{18} For any events A and B Given $P(A\cup B)=0.6$, $P(A)=P(B)$, $P(B/A)=0.8$. Then the value of $P\left[\left(A\cap\bar B\right)\cup\left(\bar A\cap B\right)\right]$ is
\begin{qoptions}
\opt{A}{$\dfrac13$}
\opt{B}{$\dfrac12$}
\opt{C}{$\dfrac14$}
\opt{D}{$\dfrac15$}
\end{qoptions}
\qn{19} If AB is a chord of the circle $x^2+y^2=25$. The tangent of A and B intersect at C. If $(2,3)$ is the mid point of AB, then the area of the quadrilateral OACB is, where O is the origin.
\begin{qoptions}
\opt{A}{$\sqrt{\dfrac{3}{13}}$ sq. unit}
\opt{B}{$50\sqrt{\dfrac{3}{13}}$ sq. unit}
\opt{C}{$100$ sq. unit}
\opt{D}{$100\sqrt3$ sq. unit}
\end{qoptions}
\qn{20} A company has 10 employees. The company has decided to form a team including atleast three employees and also excluding atleast three employees. Then the number of ways of forming the team is
\begin{qoptions}
\opt{A}{$950$}
\opt{B}{$912$}
\opt{C}{$1000$}
\opt{D}{$953$}
\end{qoptions}
\sect{Section-II}
\qn{1} If $C_1:x^2+y^2=\left(3+2\sqrt2\right)^2$ be a circle. PA and PB are pair of tangents on $C_1$ Where $P$ is any point on the director circle of $C_1$, then the radius of the smallest circle which touches $C_1$ externally and also the two tangents PA and PB is
\qn{2} The probability of getting either all heads or all tails for exactly the second time in the 3$^{\text{rd}}$ trial, if in each trial three coins are tossed, is -
\qn{3} Let the radius of the circle which passes through the point $(1,1)$ and which touches the circle $x^2+y^2+4x-6y-3=0$ at the point $(2,3)$ on it is $r$. Then the value of $r$ is
\qn{4} If the middle term in the binomial expression of $\left(x^2+\dfrac{1}{2x^2}\sin^{-1}(2x)\right)^{10}$ is equal to $\dfrac{7}{216}\left(\dfrac\pi2\right)^5$, then value of $x$ is -
\qn{5} Number of words that can be made with the letters of the word ``GENIUS'' if each word neither beginning with G nor ends with S is
\qn{6} The value of $\displaystyle\sum_{r=0}^{8}{}^8C_r\,{}^{17}C_{16-r}$ is equal to $^pC_q$, then maximum value of $(p-q)$ is
\qn{7} A regular polygon of $n$ sides has 54 diagonals. If $T$ is the number of triangles formed with exactly one side of the triangle coinciding with the side of polygon, then $\dfrac{T}{16}$ is
\qn{8} Tangents drawn from $P(1,8)$ to the circle $x^2+y^2-6x-4y-11=0$ touches the circle at the points A and B, respectively. The radius of the circle which passes through the point of intersection of circles $x^2+y^2-2x-6y+6=0$ and $x^2+y^2+2x-6y+6=0$ and intersects the circumcircle of the $\Delta PAB$ orthogonally is $R$, then $4R^2$ is equal to :-
\qn{9} The chord of contact of tangents drawn from a point on the circle $x^2+y^2=16$ to the circle $x^2+y^2=\lambda^2$ touches the circle $x^2+y^2=1$. Then value of $|\lambda|$ is equal to
\qn{10} Number of cyphers at the end of $^{100}C_{50}$ is
\end{multicols}
\end{document}
